% ch7-spectral-theorem.tex — Chapter 7: Spectral Theorem and Symmetric Matrices
% Self-adjoint operators, spectral theorem (real and complex), orthogonal
% diagonalization, normal operators, positive definiteness, SVD, applications.

\chapter{Spectral Theorem and Symmetric Matrices}\label{ch:spectral}

% ============================================================
% Motivating introduction
% ============================================================

In the preceding chapters we studied eigenvalues, eigenvectors, and
diagonalization of general linear maps.  We saw that not every matrix
is diagonalizable, and that when diagonalization succeeds the change-of-basis
matrix need not enjoy any particular geometric property.  In this chapter
we restrict our attention to a distinguished class of operators ---
\emph{self-adjoint}\index{self-adjoint operator} (symmetric or Hermitian)
operators --- for which the theory becomes strikingly clean: every
self-adjoint operator is diagonalizable, all eigenvalues are real, and
one can find an \emph{orthonormal} basis of eigenvectors.

This is the content of the \emph{Spectral Theorem},\index{Spectral Theorem}
one of the most important results in all of linear algebra.  Its applications
pervade mathematics, science, and engineering:
\begin{itemize}
  \item \textbf{Covariance matrices.}  In statistics and data science, the
    covariance matrix of a random vector is symmetric and positive
    semi\-definite; its eigenvectors are the \emph{principal components}, and
    PCA\index{PCA}\index{principal component analysis} amounts to the spectral
    decomposition of this matrix.
  \item \textbf{Hessians and optimization.}  The Hessian of a twice
    differentiable function $f\colon\R^n\to\R$ is a real symmetric matrix;
    its eigenvalues determine whether a critical point is a local minimum,
    maximum, or saddle point.
  \item \textbf{Quadratic forms.}  The classification of quadratic forms
    (ellipses, hyperbolas, etc.)\ relies on the signs of the eigenvalues of
    the associated symmetric matrix.
  \item \textbf{Vibrations and normal modes.}  The natural frequencies of a
    vibrating system are the square roots of the eigenvalues of a symmetric
    matrix, and the normal modes are its eigenvectors.
  \item \textbf{Quantum mechanics.}  Observables are Hermitian operators on a
    Hilbert space; the Spectral Theorem guarantees that measurement outcomes
    (eigenvalues) are real numbers.
\end{itemize}

\begin{center}
\begin{tikzpicture}[>=Stealth, scale=1.2]
  % --- Conic classification via eigenvalues ---
  % Ellipse
  \begin{scope}[shift={(-4,0)}]
    \draw[thick, blue!70!black] (0,0) ellipse (1.5 and 0.8);
    \draw[->, thick, red!70!black] (0,0) -- (1.5,0)
      node[below right] {\small $\lambda_1>0$};
    \draw[->, thick, red!70!black] (0,0) -- (0,0.8)
      node[above right] {\small $\lambda_2>0$};
    \node[below] at (0,-1.2) {\small Ellipse};
  \end{scope}

  % Hyperbola
  \begin{scope}[shift={(0,0)}]
    \draw[thick, blue!70!black, domain=-1.2:1.2, samples=80]
      plot ({cosh(\x)}, {0.7*sinh(\x)});
    \draw[thick, blue!70!black, domain=-1.2:1.2, samples=80]
      plot ({-cosh(\x)}, {0.7*sinh(\x)});
    \draw[->, thick, red!70!black] (0,0) -- (1.3,0)
      node[below right] {\small $\lambda_1>0$};
    \draw[->, thick, green!60!black] (0,0) -- (0,0.9)
      node[above right] {\small $\lambda_2<0$};
    \node[below] at (0,-1.2) {\small Hyperbola};
  \end{scope}

  % Parabola (degenerate: one eigenvalue zero)
  \begin{scope}[shift={(4,0)}]
    \draw[thick, blue!70!black, domain=-1:1, samples=60]
      plot ({1.2*\x*\x}, {\x});
    \draw[->, thick, red!70!black] (0,0) -- (1.3,0)
      node[below right] {\small $\lambda_1>0$};
    \draw[->, thick, gray] (0,0) -- (0,0.9)
      node[above right] {\small $\lambda_2=0$};
    \node[below] at (0.5,-1.2) {\small Parabola};
  \end{scope}
\end{tikzpicture}
\captionof{figure}{Classification of conics by the eigenvalues of the
  associated symmetric matrix.  All eigenvalues positive: ellipse;
  mixed signs: hyperbola; one eigenvalue zero: parabola.}
\label{fig:conic-classification}
\end{center}

Throughout this chapter, $E$ denotes a finite-dimensional inner product
space over $\K = \R$ or~$\C$, with inner product $\inner{{\cdot},{\ \cdot}}$
and induced norm~$\norm{\cdot}$.  We write $\dim E = n$.  Matrices are
identified with their corresponding endomorphisms via a chosen basis
whenever convenient.


% ============================================================
\section{Self-adjoint (symmetric and Hermitian) operators}%
\label{sec:self-adjoint}
% ============================================================

\begin{definition}{Adjoint of a linear map}{def:adjoint-ch7}
\index{adjoint operator}
Let $(E,\inner{\cdot,\cdot})$ be a finite-dimensional inner product
space and $f\in\End(E)$.  The \emph{adjoint} of~$f$ is the unique
linear map $f^*\in\End(E)$ satisfying
\[
  \inner{f(u),\, v} \;=\; \inner{u,\, f^*(v)}
  \qquad\text{for all } u,v\in E.
\]
\end{definition}

\begin{remark}{Existence and uniqueness of the adjoint}{rem:adjoint-exists}
Existence and uniqueness follow from the Riesz representation theorem
(or can be verified directly using an orthonormal basis).  If $\cB$ is
an orthonormal basis and $A = \Mat_\cB(f)$, then
$\Mat_\cB(f^*) = \hermit{A}$, the conjugate transpose.  Over~$\R$
this reduces to $\Mat_\cB(f^*) = \transpose{A}$.
\end{remark}

\begin{proposition}{Properties of the adjoint}{prop:adjoint-properties-ch7}
\index{adjoint operator!properties}
For $f,g\in\End(E)$ and $\alpha\in\K$:
\begin{enumerate}[label=\textup{(\roman*)}]
  \item $(f+g)^* = f^*+g^*$.
  \item $(\alpha f)^* = \bar\alpha\, f^*$.
  \item $(f\circ g)^* = g^*\circ f^*$.
  \item $(f^*)^* = f$.
  \item $f$ is invertible if and only if $f^*$ is, and then
    $(f^*)^{-1}=(f^{-1})^*$.
\end{enumerate}
\end{proposition}

\begin{proof}
All properties follow from the definition and the sesquilinearity
(or bilinearity over~$\R$) of the inner product.  We verify~(iii):
for all $u,v\in E$,
\[
  \inner{(f\circ g)(u),\,v}
  = \inner{f(g(u)),\,v}
  = \inner{g(u),\,f^*(v)}
  = \inner{u,\,g^*(f^*(v))}
  = \inner{u,\,(g^*\circ f^*)(v)},
\]
so $(f\circ g)^* = g^*\circ f^*$ by uniqueness of the adjoint.
The remaining items are proved similarly.
\end{proof}

\begin{definition}{Self-adjoint operator}{def:self-adjoint}
\index{self-adjoint operator}\index{symmetric matrix}\index{Hermitian matrix}
An endomorphism $f\in\End(E)$ is \emph{self-adjoint} if $f^* = f$,
i.e.\ $\inner{f(u),v} = \inner{u,f(v)}$ for all $u,v\in E$.

In matrix terms with respect to an orthonormal basis:
\begin{itemize}
  \item Over $\R$: $A$ is \emph{symmetric},\index{symmetric matrix}
    i.e.\ $\transpose{A} = A$.
  \item Over $\C$: $A$ is \emph{Hermitian},\index{Hermitian matrix}
    i.e.\ $\hermit{A} = A$.
\end{itemize}

We denote the real vector space of $n\times n$ real symmetric matrices by
$\Sym_n(\R)\deff\setcond{A\in\Mat_n(\R)}{\transpose{A}=A}$.
\end{definition}

\begin{example}{Symmetric matrices}{ex:symmetric-examples}
The matrix
$A = \begin{pmatrix} 2 & -1 \\ -1 & 3 \end{pmatrix}$
is symmetric.  The matrix
$B = \begin{pmatrix} 1 & 2+i \\ 2-i & 3 \end{pmatrix}$
is Hermitian (the diagonal entries are real and the off-diagonal entries
are complex conjugates).
\end{example}

\begin{proposition}{Dimension of $\Sym_n(\R)$}{prop:dim-sym}
$\dim\Sym_n(\R) = \dfrac{n(n+1)}{2}$.
\end{proposition}

\begin{proof}
A symmetric matrix is determined by its entries on and above the diagonal.
There are $n$ diagonal entries and $\binom{n}{2}$ entries strictly above
the diagonal, giving $n + \binom{n}{2} = \frac{n(n+1)}{2}$.
\end{proof}


% ============================================================
\section{Eigenvalues of self-adjoint operators are real}%
\label{sec:eigenvalues-real}
% ============================================================

\begin{theorem}{Eigenvalues of self-adjoint operators are real}%
{thm:eigenvalues-real}
\index{eigenvalue!of self-adjoint operator}\index{self-adjoint operator!eigenvalues}
Let $f$ be a self-adjoint operator on an inner product space~$E$
(over $\R$ or~$\C$).  Then every eigenvalue of~$f$ is real.
\end{theorem}

\begin{proof}
Let $\lambda$ be an eigenvalue and $v\neq\vzero$ an eigenvector:
$f(v)=\lambda v$.  Then
\[
  \lambda\,\inner{v,v}
  = \inner{\lambda v,\,v}
  = \inner{f(v),\,v}
  = \inner{v,\,f(v)}
  = \inner{v,\,\lambda v}
  = \bar\lambda\,\inner{v,v}.
\]
Since $v\neq\vzero$, we have $\inner{v,v}>0$, so dividing yields
$\lambda = \bar\lambda$, which means $\lambda\in\R$.
\end{proof}

\begin{remark}{Real eigenvalues over $\R$}{rem:real-eigenvalues-R}
Over~$\R$, one must be slightly careful: a real symmetric matrix might
a priori have no eigenvalues in~$\R$ (since the characteristic polynomial
need not split).  However, the Spectral Theorem will establish that it
does.  The above proof works directly if we allow eigenvalues in~$\C$
and use the standard Hermitian inner product on~$\C^n$.  Concretely:
if $A\in\Sym_n(\R)$ and $\lambda\in\C$ satisfies $Av = \lambda v$ for
some $v\in\C^n\setminus\set{\vzero}$, then $\lambda\in\R$.
\end{remark}


% ============================================================
\section{Orthogonality of eigenvectors for distinct eigenvalues}%
\label{sec:eigenvectors-orthogonal}
% ============================================================

\begin{theorem}{Eigenvectors for distinct eigenvalues are orthogonal}%
{thm:eigenvectors-orthogonal}
\index{eigenvector!orthogonality}\index{self-adjoint operator!orthogonal eigenvectors}
Let $f$ be a self-adjoint operator on~$E$, and let $v_1,v_2$ be
eigenvectors of~$f$ associated with distinct eigenvalues
$\lambda_1\neq\lambda_2$.  Then $\inner{v_1,v_2}=0$.
\end{theorem}

\begin{proof}
We compute $\inner{f(v_1),v_2}$ in two ways:
\[
  \inner{f(v_1),\,v_2}
  = \inner{\lambda_1 v_1,\,v_2}
  = \lambda_1\,\inner{v_1,v_2},
\]
and, using the self-adjointness of~$f$,
\[
  \inner{f(v_1),\,v_2}
  = \inner{v_1,\,f(v_2)}
  = \inner{v_1,\,\lambda_2 v_2}
  = \bar\lambda_2\,\inner{v_1,v_2}
  = \lambda_2\,\inner{v_1,v_2},
\]
where the last equality uses the fact that $\lambda_2\in\R$
(
ef{thm:thm:eigenvalues-real}).  Subtracting:
\[
  (\lambda_1 - \lambda_2)\,\inner{v_1,v_2} = 0.
\]
Since $\lambda_1\neq\lambda_2$, we conclude $\inner{v_1,v_2}=0$.
\end{proof}

\begin{corollary}{Eigenspaces of a self-adjoint operator are mutually
orthogonal}{cor:eigenspaces-orthogonal}
\index{eigenspace!orthogonality}
If $f$ is self-adjoint with distinct eigenvalues
$\lambda_1,\ldots,\lambda_r$, then
\[
  E_{\lambda_i} \perp E_{\lambda_j}
  \qquad\text{for all } i\neq j,
\]
where $E_{\lambda_k} = \Ker(f - \lambda_k\Id)$ is the eigenspace for
$\lambda_k$.
\end{corollary}

\begin{proof}
Every vector in $E_{\lambda_i}$ is an eigenvector for~$\lambda_i$
(or zero), and likewise for~$E_{\lambda_j}$.  The result follows from

ef{thm:thm:eigenvectors-orthogonal} applied to each pair of nonzero vectors.
\end{proof}


% ============================================================
\section{The Spectral Theorem for real symmetric matrices}%
\label{sec:spectral-real}
% ============================================================

We now prove the central result of this chapter.

\begin{lemma}{Invariance of the orthogonal complement}%
{lem:orth-complement-invariant}
\index{orthogonal complement!invariance under self-adjoint maps}
Let $f$ be a self-adjoint operator on~$E$ and let $W\subseteq E$ be a
subspace invariant under~$f$ (i.e.\ $f(W)\subseteq W$).  Then the
orthogonal complement $W^\perp$ is also invariant under~$f$.
\end{lemma}

\begin{proof}
Let $v\in W^\perp$.  We must show $f(v)\in W^\perp$, i.e.\
$\inner{f(v),w}=0$ for all $w\in W$.  Since $f$ is self-adjoint,
\[
  \inner{f(v),\,w} = \inner{v,\,f(w)}.
\]
Now $w\in W$ and $f(W)\subseteq W$, so $f(w)\in W$.  Since
$v\in W^\perp$, we get $\inner{v,f(w)}=0$.  Hence $\inner{f(v),w}=0$.
\end{proof}

\begin{lemma}{A self-adjoint operator over $\R$ has a real eigenvalue}%
{lem:self-adjoint-has-eigenvalue}
\index{self-adjoint operator!existence of eigenvalue}
Let $A\in\Sym_n(\R)$.  Then $A$ has at least one (real) eigenvalue.
\end{lemma}

\begin{proof}
The characteristic polynomial $\chi_A(\lambda) = \det(A - \lambda I_n)$
has degree~$n$ with real coefficients.  Over~$\C$ it has a root
$\lambda_0\in\C$, and by 
ef{thm:thm:eigenvalues-real} (applied in the
Hermitian setting) $\lambda_0$ is real.  Hence $A$ has a real eigenvalue.
\end{proof}

\begin{theorem}{Spectral Theorem for real symmetric matrices}%
{thm:spectral-real}
\index{Spectral Theorem!real symmetric matrices}
Let $A\in\Sym_n(\R)$ be a real symmetric matrix.  Then:
\begin{enumerate}[label=\textup{(\roman*)}]
  \item All eigenvalues of~$A$ are real.
  \item There exists an orthogonal matrix $P\in\Orth(n)$
    (i.e.\ $\transpose{P}P = I_n$) such that
    \[
      \transpose{P}\,A\,P = \diag(\lambda_1,\ldots,\lambda_n),
    \]
    where $\lambda_1,\ldots,\lambda_n$ are the eigenvalues of~$A$
    (counted with multiplicity).
  \item Equivalently, $\R^n$ has an orthonormal basis of eigenvectors
    of~$A$.
\end{enumerate}
\end{theorem}

\begin{proof}
We prove by strong induction on~$n$.

\medskip\noindent\textbf{Base case ($n=1$).}\quad
A $1\times 1$ matrix $A=(a)$ is already diagonal with
eigenvalue~$a\in\R$, and $P=(1)\in\Orth(1)$.

\medskip\noindent\textbf{Inductive step.}\quad
Assume the theorem holds for all real symmetric matrices of size less
than~$n$.  Let $A\in\Sym_n(\R)$.

\medskip\noindent\emph{Step~1: Find a real eigenvalue and unit eigenvector.}\quad
By 
ef{lem:lem:self-adjoint-has-eigenvalue}, $A$ has a real eigenvalue
$\lambda_1\in\R$.  Let $v_1$ be an eigenvector for~$\lambda_1$ with
$\norm{v_1}=1$.

\medskip\noindent\emph{Step~2: Set up the orthogonal complement.}\quad
Let $W = \Span(v_1)$.  Since $A v_1 = \lambda_1 v_1$, the subspace~$W$
is $A$-invariant.  By 
ef{lem:lem:orth-complement-invariant}, $W^\perp$
is also $A$-invariant, and $\dim W^\perp = n-1$.

\medskip\noindent\emph{Step~3: Restrict and apply the induction hypothesis.}\quad
Let $g = f\restrict{W^\perp}$, where $f$ is the linear map
$x\mapsto Ax$.  Since $f(W^\perp)\subseteq W^\perp$, the map~$g$ is a
well-defined endomorphism of~$W^\perp$.  Moreover, $g$ is self-adjoint
on~$W^\perp$ (the inner product is the restriction of the standard one,
and the self-adjointness condition is inherited).

Choose any orthonormal basis of~$W^\perp$ and let $B$ be the matrix
of~$g$ in this basis.  Then $B\in\Sym_{n-1}(\R)$.  By the induction
hypothesis, there exists an orthonormal basis
$(v_2,\ldots,v_n)$ of~$W^\perp$ consisting of eigenvectors of~$g$
(equivalently, of~$f$), with real eigenvalues
$\lambda_2,\ldots,\lambda_n$.

\medskip\noindent\emph{Step~4: Assemble the orthonormal eigenbasis.}\quad
The family $(v_1,v_2,\ldots,v_n)$ is orthonormal:
\begin{itemize}
  \item $v_1\perp v_j$ for $j\geq 2$, since $v_j\in W^\perp$;
  \item $(v_2,\ldots,v_n)$ is orthonormal by the induction hypothesis;
  \item $\norm{v_1}=1$ by construction.
\end{itemize}
Each $v_j$ is an eigenvector of~$A$ with eigenvalue~$\lambda_j$.
Let $P = \begin{pmatrix} v_1 & v_2 & \cdots & v_n \end{pmatrix}$.
Then $P\in\Orth(n)$ and $\transpose{P}AP = \diag(\lambda_1,\ldots,\lambda_n)$.
\end{proof}

\begin{remark}{Spectral decomposition form}{rem:spectral-decomposition}
\index{spectral decomposition}
The Spectral Theorem can be rephrased as follows.  If $A\in\Sym_n(\R)$
has distinct eigenvalues $\mu_1,\ldots,\mu_r$ with orthogonal projections
$P_i$ onto the eigenspace $E_{\mu_i}$, then
\[
  A = \mu_1 P_1 + \mu_2 P_2 + \cdots + \mu_r P_r,
\]
where $P_1 + \cdots + P_r = I_n$ and $P_i P_j = 0$ for $i\neq j$.
This is called the \emph{spectral decomposition}\index{spectral decomposition}
of~$A$.
\end{remark}


% ============================================================
\section{The Spectral Theorem for Hermitian matrices}%
\label{sec:spectral-hermitian}
% ============================================================

The real Spectral Theorem extends naturally to the complex setting.

\begin{theorem}{Spectral Theorem for Hermitian matrices}%
{thm:spectral-hermitian}
\index{Spectral Theorem!Hermitian matrices}
Let $A\in\Mat_n(\C)$ be Hermitian, i.e.\ $\hermit{A}=A$.  Then:
\begin{enumerate}[label=\textup{(\roman*)}]
  \item All eigenvalues of~$A$ are real.
  \item There exists a unitary matrix $U\in\Mat_n(\C)$
    (i.e.\ $\hermit{U}U = I_n$) such that
    \[
      \hermit{U}\,A\,U = \diag(\lambda_1,\ldots,\lambda_n),
    \]
    with $\lambda_1,\ldots,\lambda_n\in\R$.
  \item Equivalently, $\C^n$ has an orthonormal basis of eigenvectors
    of~$A$.
\end{enumerate}
\end{theorem}

\begin{proof}
The proof is identical to the real case: over~$\C$ the characteristic
polynomial always splits, so the base step of the induction is
guaranteed.  The self-adjointness condition $\hermit{A}=A$ ensures that
eigenvalues are real (
ef{thm:thm:eigenvalues-real}), eigenvectors for
distinct eigenvalues are orthogonal
(
ef{thm:thm:eigenvectors-orthogonal}), and the orthogonal complement
of an invariant subspace is invariant
(
ef{lem:lem:orth-complement-invariant}).  The inductive argument carries
over verbatim, with ``orthogonal matrix'' replaced by ``unitary matrix''
and ``transpose'' replaced by ``conjugate transpose.''
\end{proof}


% ============================================================
\section{Orthogonal diagonalization: procedure and examples}%
\label{sec:orth-diag}
% ============================================================

The Spectral Theorem provides both a theoretical guarantee and a
practical algorithm for orthogonally diagonalizing a symmetric matrix.

\begin{remark}{Orthogonal diagonalization procedure}{rem:orth-diag-procedure}
\index{orthogonal diagonalization!procedure}
Given $A\in\Sym_n(\R)$:
\begin{enumerate}
  \item Compute the characteristic polynomial
    $\chi_A(\lambda) = \det(A - \lambda I_n)$.
  \item Find all eigenvalues $\lambda_1,\ldots,\lambda_r$ (they are real).
  \item For each~$\lambda_i$, solve $(A - \lambda_i I_n)X = \vzero$ to
    find a basis of the eigenspace $E_{\lambda_i}$.
  \item Apply the Gram--Schmidt process within each eigenspace to obtain
    an orthonormal basis of~$E_{\lambda_i}$.
    (Eigenvectors from \emph{different} eigenspaces are automatically
    orthogonal, so Gram--Schmidt is only needed within a single
    eigenspace when $\dim E_{\lambda_i}\geq 2$.)
  \item Form $P$ by placing the orthonormal eigenvectors as columns.
    Then $P\in\Orth(n)$ and $\transpose{P}AP = \diag(\lambda_1,\ldots,\lambda_n)$.
\end{enumerate}
\end{remark}

\begin{example}{Orthogonal diagonalization of a $2\times 2$ matrix}%
{ex:orth-diag-2x2}
Diagonalize $A = \begin{pmatrix} 5 & 2 \\ 2 & 2 \end{pmatrix}$
orthogonally.

\emph{Step~1.} $\chi_A(\lambda) = (5-\lambda)(2-\lambda)-4
= \lambda^2-7\lambda+6 = (\lambda-1)(\lambda-6)$.

\emph{Step~2.} Eigenvalues: $\lambda_1=1$, $\lambda_2=6$.

\emph{Step~3.}
For $\lambda_1=1$: $(A-I_2)X=\vzero$ gives
$\begin{pmatrix}4&2\\2&1\end{pmatrix}X=\vzero$, so $2x+y=0$.
Eigenvector: $v_1 = \begin{psmallmatrix}1\\-2\end{psmallmatrix}$.

For $\lambda_2=6$: $(A-6I_2)X=\vzero$ gives
$\begin{pmatrix}-1&2\\2&-4\end{pmatrix}X=\vzero$, so $x=2y$.
Eigenvector: $v_2 = \begin{psmallmatrix}2\\1\end{psmallmatrix}$.

\emph{Step~4.} Normalise:
$\hat v_1 = \frac{1}{\sqrt{5}}\begin{psmallmatrix}1\\-2\end{psmallmatrix}$,
$\hat v_2 = \frac{1}{\sqrt{5}}\begin{psmallmatrix}2\\1\end{psmallmatrix}$.

\emph{Step~5.}
$P = \frac{1}{\sqrt{5}}\begin{pmatrix}1&2\\-2&1\end{pmatrix}$
satisfies $\transpose{P}AP = \begin{pmatrix}1&0\\0&6\end{pmatrix}$.
\end{example}

\begin{example}{Orthogonal diagonalization of a $3\times 3$ matrix}%
{ex:orth-diag-3x3}
Let
$A = \begin{pmatrix} 2 & 1 & 1 \\ 1 & 2 & 1 \\ 1 & 1 & 2 \end{pmatrix}$.
One verifies that $\transpose{A}=A$.

\emph{Characteristic polynomial:}
$\chi_A(\lambda) = -(\lambda-4)(\lambda-1)^2$.

\emph{Eigenvalues:} $\lambda_1=4$ (multiplicity~$1$),
$\lambda_2=1$ (multiplicity~$2$).

\emph{For $\lambda_1=4$:}
$(A-4I_3)X=\vzero$ gives $x_1=x_2=x_3$, so
$E_4 = \Span\!\left(\begin{psmallmatrix}1\\1\\1\end{psmallmatrix}\right)$.
Normalise: $u_1 = \frac{1}{\sqrt{3}}\begin{psmallmatrix}1\\1\\1\end{psmallmatrix}$.

\emph{For $\lambda_2=1$:}
$(A-I_3)X=\vzero$ gives $x_1+x_2+x_3=0$, so
$E_1 = \Span\!\left(
  \begin{psmallmatrix}1\\-1\\0\end{psmallmatrix},\,
  \begin{psmallmatrix}1\\0\\-1\end{psmallmatrix}
\right)$.
Apply Gram--Schmidt:
$w_1 = \begin{psmallmatrix}1\\-1\\0\end{psmallmatrix}$,
$w_2 = \begin{psmallmatrix}1\\0\\-1\end{psmallmatrix}
  - \frac{\inner{w_1,\begin{psmallmatrix}1\\0\\-1\end{psmallmatrix}}}{\inner{w_1,w_1}}\,w_1
= \begin{psmallmatrix}1/2\\1/2\\-1\end{psmallmatrix}$.

Normalise:
$u_2 = \frac{1}{\sqrt{2}}\begin{psmallmatrix}1\\-1\\0\end{psmallmatrix}$,
$u_3 = \frac{1}{\sqrt{6}}\begin{psmallmatrix}1\\1\\-2\end{psmallmatrix}$.

The orthogonal matrix is
$P = \begin{pmatrix}
\frac{1}{\sqrt{3}} & \frac{1}{\sqrt{2}} & \frac{1}{\sqrt{6}} \\[4pt]
\frac{1}{\sqrt{3}} & -\frac{1}{\sqrt{2}} & \frac{1}{\sqrt{6}} \\[4pt]
\frac{1}{\sqrt{3}} & 0 & -\frac{2}{\sqrt{6}}
\end{pmatrix}$,
and $\transpose{P}AP = \diag(4,1,1)$.
\end{example}


% ============================================================
\section{Normal operators}\label{sec:normal}
% ============================================================

Self-adjoint operators are a special case of a broader class.

\begin{definition}{Normal operator}{def:normal}
\index{normal operator}\index{normal matrix}
An operator $f\in\End(E)$ is \emph{normal} if it commutes with its
adjoint:
\[
  f\circ f^* = f^*\circ f.
\]
A matrix $A\in\Mat_n(\C)$ is \emph{normal} if $A\hermit{A} = \hermit{A}A$.
\end{definition}

\begin{example}{Classes of normal operators}{ex:normal-classes}
The following are all normal:
\begin{itemize}
  \item Self-adjoint (symmetric / Hermitian) operators: $f^*=f$.
  \item Skew-adjoint operators: $f^*=-f$ (equivalently,
    $\transpose{A}=-A$ over~$\R$ or $\hermit{A}=-A$ over~$\C$).
  \item Orthogonal / unitary operators: $f^*=\inv{f}$.
\end{itemize}
\end{example}

\begin{proposition}{Characterization of normality}{prop:normal-char}
Let $f\in\End(E)$ where $E$ is a finite-dimensional inner product space
over~$\C$.  The following are equivalent:
\begin{enumerate}[label=\textup{(\roman*)}]
  \item $f$ is normal.
  \item $\norm{f(v)} = \norm{f^*(v)}$ for all $v\in E$.
  \item If $f(v)=\lambda v$, then $f^*(v)=\bar\lambda v$.
    (That is, $v$ is an eigenvector of~$f$ if and only if it is an
    eigenvector of~$f^*$, with conjugate eigenvalue.)
\end{enumerate}
\end{proposition}

\begin{proof}
(i)$\Rightarrow$(ii): $\norm{f(v)}^2 = \inner{f(v),f(v)}
= \inner{f^*f(v),v} = \inner{ff^*(v),v} = \inner{f^*(v),f^*(v)}
= \norm{f^*(v)}^2$.

(ii)$\Rightarrow$(iii): Assume $f(v)=\lambda v$.  Apply~(ii) to
$f-\lambda\Id$ (which is also normal, as one verifies):
$\norm{(f-\lambda\Id)(v)}=\norm{(f-\lambda\Id)^*(v)}
=\norm{(f^*-\bar\lambda\Id)(v)}$.
The left side is~$0$, so $f^*(v)=\bar\lambda v$.

(iii)$\Rightarrow$(i): We show $f^*f = ff^*$ by checking on an
eigenbasis.  Over~$\C$, every normal operator is unitarily
diagonalizable (see the next theorem), so this direction also follows.
Alternatively, one can verify
$\inner{(f^*f-ff^*)(v),v}=0$ for all~$v$ using~(iii) and polarization.
\end{proof}

\begin{theorem}{Spectral Theorem for normal operators over $\C$}%
{thm:spectral-normal}
\index{Spectral Theorem!normal operators}
Let $A\in\Mat_n(\C)$ be normal.  Then there exists a unitary matrix
$U$ such that
\[
  \hermit{U}\,A\,U = \diag(\lambda_1,\ldots,\lambda_n),
\]
where $\lambda_1,\ldots,\lambda_n\in\C$ are the eigenvalues of~$A$.

Conversely, every unitarily diagonalizable matrix is normal.
\end{theorem}

\begin{proof}
\emph{Forward direction.}  By induction on~$n$.  The case $n=1$ is
trivial.  For the inductive step, since $\C$ is algebraically closed,
$A$ has an eigenvalue $\lambda_1$ with unit eigenvector~$v_1$.
Let $W=\Span(v_1)$.

\emph{Claim:} $W^\perp$ is $A$-invariant.  Let $u\in W^\perp$.  We need
$Au\perp v_1$, i.e.\ $\inner{Au,v_1}=0$.  Now
$\inner{Au,v_1} = \inner{u,\hermit{A}v_1} = \inner{u,\bar\lambda_1 v_1}
= \lambda_1\inner{u,v_1} = 0$,
where we used 
ef{prop:prop:normal-char}~(iii): since $Av_1=\lambda_1 v_1$,
we have $\hermit{A}v_1=\bar\lambda_1 v_1$.

The restriction $A\restrict{W^\perp}$ is normal on the $(n-1)$-dimensional
space~$W^\perp$.  By the induction hypothesis it is unitarily
diagonalizable.  Assembling the eigenvectors gives a unitary matrix~$U$
diagonalizing~$A$.

\emph{Converse.}  If $A = U D \hermit{U}$ with $D$ diagonal and $U$
unitary, then $\hermit{A} = U\bar D\hermit{U}$.  Since diagonal
matrices commute, $A\hermit{A} = UD\bar D\hermit{U}
= U\bar D D\hermit{U} = \hermit{A}A$.
\end{proof}

\begin{remark}{Normal operators over $\R$}{rem:normal-real}
\index{normal operator!over $\R$}
Over~$\R$, normal operators need not be orthogonally diagonalizable.
For instance, the rotation matrix
$R_\theta = \begin{psmallmatrix}\cos\theta & -\sin\theta \\
\sin\theta & \cos\theta\end{psmallmatrix}$
is normal (in fact orthogonal) but has no real eigenvalues when
$\theta\notin\set{0,\pi}$.  The real spectral theorem for normal
operators states that such a matrix can be brought to a block-diagonal
form with $1\times 1$ real blocks and $2\times 2$ rotation blocks.
\end{remark}


% ============================================================
\section{Positive definite and positive semidefinite matrices}%
\label{sec:positive-definite}
% ============================================================

\begin{definition}{Positive definite and positive semidefinite matrices}%
{def:positive-definite}
\index{positive definite matrix}\index{positive semidefinite matrix}
A real symmetric matrix $A\in\Sym_n(\R)$ is:
\begin{itemize}
  \item \emph{positive definite}\index{positive definite matrix}
    if $\transpose{x}Ax > 0$ for all $x\in\R^n\setminus\set{\vzero}$;
  \item \emph{positive semidefinite}\index{positive semidefinite matrix}
    if $\transpose{x}Ax \geq 0$ for all $x\in\R^n$;
  \item \emph{negative definite} if $-A$ is positive definite;
  \item \emph{negative semidefinite} if $-A$ is positive semidefinite;
  \item \emph{indefinite} if $\transpose{x}Ax$ takes both positive and
    negative values.
\end{itemize}
\end{definition}

\begin{theorem}{Characterizations of positive definiteness}%
{thm:pd-characterizations}
\index{positive definite matrix!characterizations}
For $A\in\Sym_n(\R)$ the following are equivalent:
\begin{enumerate}[label=\textup{(\roman*)}]
  \item $A$ is positive definite.
  \item All eigenvalues of~$A$ are strictly positive.
  \item All leading principal minors of~$A$ are strictly positive
    (\emph{Sylvester's criterion}\index{Sylvester's criterion}).
  \item There exists an invertible matrix $L\in\Mat_n(\R)$ such that
    $A = \transpose{L}\,L$.
  \item $A$ admits a \emph{Cholesky decomposition}\index{Cholesky decomposition}:
    there exists a unique lower triangular matrix $L$ with strictly positive
    diagonal entries such that $A = L\transpose{L}$.
\end{enumerate}
\end{theorem}

\begin{proof}
(i)$\Rightarrow$(ii): If $\lambda$ is an eigenvalue with eigenvector~$v$,
then $0 < \transpose{v}Av = \lambda\transpose{v}v = \lambda\norm{v}^2$,
so $\lambda > 0$.

(ii)$\Rightarrow$(i): By the Spectral Theorem, $A = P\,D\,\transpose{P}$
with $P$ orthogonal and $D = \diag(\lambda_1,\ldots,\lambda_n)$.  For
$x\neq\vzero$, let $y = \transpose{P}x \neq \vzero$.  Then
$\transpose{x}Ax = \transpose{y}Dy = \sum_{i=1}^n \lambda_i y_i^2 > 0$
since all $\lambda_i>0$ and $y\neq\vzero$.

(i)$\Rightarrow$(iii) (\emph{Sylvester's criterion}): The $k$-th
leading principal minor $\Delta_k$ equals the determinant of the
$k\times k$ upper-left submatrix~$A_k$.  But $A_k$ is itself symmetric,
and for any $x\in\R^k\setminus\set{\vzero}$ we have
$\transpose{x}A_k x = \transpose{\tilde x}A\tilde x > 0$, where
$\tilde x = (x_1,\ldots,x_k,0,\ldots,0)^T$.  Hence $A_k$ is positive
definite, so its eigenvalues are positive, so
$\Delta_k = \det(A_k) > 0$.

(iii)$\Rightarrow$(i): By induction on~$n$.  For $n=1$, $\Delta_1 = a_{11} > 0$,
so $A=(a_{11})$ is positive definite.  For the inductive step, write
$A = \begin{psmallmatrix} A_{n-1} & b \\ \transpose{b} & a_{nn}\end{psmallmatrix}$.
By induction, $A_{n-1}$ is positive definite.  Performing one step of block
Gaussian elimination (valid since $\det(A_{n-1})=\Delta_{n-1}>0$):
\[
  A = \begin{pmatrix} I & 0 \\ \transpose{b}\inv{A_{n-1}} & 1\end{pmatrix}
  \begin{pmatrix} A_{n-1} & 0 \\ 0 & s\end{pmatrix}
  \begin{pmatrix} I & \inv{A_{n-1}}b \\ 0 & 1\end{pmatrix},
\]
where $s = a_{nn} - \transpose{b}\inv{A_{n-1}}b
= \det(A)/\det(A_{n-1}) = \Delta_n/\Delta_{n-1} > 0$.
This expresses $A = \transpose{M}
\begin{psmallmatrix}A_{n-1}&0\\0&s\end{psmallmatrix} M$ with $M$ invertible.
For $x\neq\vzero$, $\transpose{x}Ax = \transpose{(Mx)}
\begin{psmallmatrix}A_{n-1}&0\\0&s\end{psmallmatrix}(Mx) > 0$
since the block-diagonal matrix is positive definite.

(ii)$\Rightarrow$(iv): By the Spectral Theorem, $A = PD\transpose{P}$,
so $A = PD^{1/2}\transpose{(PD^{1/2})} = \transpose{L}L$ where
$L = D^{1/2}\transpose{P}$.

(iv)$\Rightarrow$(i): $\transpose{x}Ax = \transpose{x}\transpose{L}Lx
= \norm{Lx}^2 > 0$ since $L$ is invertible and $x\neq\vzero$.

The existence and uniqueness of the Cholesky factorization~(v) is proved
by an explicit recursive construction using the positive leading
principal minors (we omit the details).
\end{proof}

\begin{example}{Testing positive definiteness}{ex:pd-test}
Consider $A = \begin{pmatrix}2 & -1 \\ -1 & 2\end{pmatrix}$.

\emph{Via eigenvalues:} $\chi_A(\lambda)=(\lambda-1)(\lambda-3)$, so
$\lambda_1=1>0$ and $\lambda_2=3>0$.  Hence $A$ is positive definite.

\emph{Via Sylvester's criterion:} $\Delta_1=2>0$, $\Delta_2=\det(A)=3>0$.
Positive definite.

\emph{Cholesky decomposition:}
$A = \begin{pmatrix}\sqrt{2}&0\\-\frac{1}{\sqrt{2}}&\frac{\sqrt{3}}{\sqrt{2}}\end{pmatrix}
\begin{pmatrix}\sqrt{2}&-\frac{1}{\sqrt{2}}\\0&\frac{\sqrt{3}}{\sqrt{2}}\end{pmatrix}
= L\transpose{L}$.
\end{example}

\begin{proposition}{Characterization of positive semidefiniteness}%
{prop:psd-characterizations}
\index{positive semidefinite matrix!characterizations}
For $A\in\Sym_n(\R)$ the following are equivalent:
\begin{enumerate}[label=\textup{(\roman*)}]
  \item $A$ is positive semidefinite.
  \item All eigenvalues of~$A$ are nonnegative.
  \item There exists a matrix $B\in\Mat_n(\R)$ such that $A=\transpose{B}B$.
  \item All principal minors of~$A$ are nonnegative.
\end{enumerate}
\end{proposition}

\begin{proof}
The proofs are analogous to 
ef{thm:thm:pd-characterizations}, with
strict inequalities relaxed to non-strict ones.  For (ii)$\Rightarrow$(iii),
use the spectral decomposition $A = PD\transpose{P}$ and set
$B = D^{1/2}\transpose{P}$ (noting $D^{1/2}$ is well-defined since all
diagonal entries are nonnegative).
\end{proof}


% ============================================================
\section{Singular Value Decomposition}\label{sec:svd}
% ============================================================

The Spectral Theorem applies to symmetric (or Hermitian) matrices.  The
\emph{Singular Value Decomposition} (SVD)\index{SVD}\index{singular value decomposition}
extends the idea of diagonalization to \emph{arbitrary} matrices,
including rectangular ones.

\begin{definition}{Singular values}{def:singular-values}
\index{singular value}
Let $A\in\Mat_{m\times n}(\R)$.  The \emph{singular values} of~$A$
are the nonnegative square roots of the eigenvalues of the positive
semidefinite matrix $\transpose{A}A\in\Sym_n(\R)$:
\[
  \sigma_i \deff \sqrt{\lambda_i(\transpose{A}A)},
  \qquad i = 1,\ldots,n,
\]
ordered so that $\sigma_1\geq\sigma_2\geq\cdots\geq\sigma_n\geq 0$.
\end{definition}

\begin{theorem}{Singular Value Decomposition}{thm:svd}
\index{SVD!theorem}\index{singular value decomposition!theorem}
Let $A\in\Mat_{m\times n}(\R)$ with $\rank(A)=r$.  Then there exist:
\begin{itemize}
  \item an orthogonal matrix $U\in\Orth(m)$ (left singular vectors),
  \item an orthogonal matrix $V\in\Orth(n)$ (right singular vectors),
  \item a ``diagonal'' matrix $\Sigma\in\Mat_{m\times n}(\R)$ with
    $\Sigma_{ii}=\sigma_i$ for $i=1,\ldots,\min(m,n)$ and all other
    entries zero,
\end{itemize}
such that
\[
  A = U\,\Sigma\,\transpose{V}.
\]
The nonzero singular values $\sigma_1\geq\cdots\geq\sigma_r>0$ are
uniquely determined.  Moreover,
\[
  A = \sum_{i=1}^{r} \sigma_i\, u_i\,\transpose{v_i},
\]
where $u_i$ and~$v_i$ are the $i$-th columns of~$U$ and~$V$.
\end{theorem}

\begin{proof}[Proof sketch]
The matrix $\transpose{A}A\in\Sym_n(\R)$ is positive semidefinite
(
ef{prop:prop:psd-characterizations}).  By the Spectral Theorem, there
exists an orthogonal matrix $V$ such that
$\transpose{V}\transpose{A}A\,V = \diag(\sigma_1^2,\ldots,\sigma_n^2)$
with $\sigma_1\geq\cdots\geq\sigma_r>0=\sigma_{r+1}=\cdots=\sigma_n$.

Define $u_i = \frac{1}{\sigma_i}Av_i$ for $i=1,\ldots,r$.  These are
orthonormal:
\[
  \inner{u_i,u_j}
  = \frac{1}{\sigma_i\sigma_j}\inner{Av_i,Av_j}
  = \frac{1}{\sigma_i\sigma_j}\transpose{v_i}\transpose{A}Av_j
  = \frac{\sigma_j^2}{\sigma_i\sigma_j}\delta_{ij}
  = \delta_{ij}.
\]
Extend $(u_1,\ldots,u_r)$ to an orthonormal basis
$(u_1,\ldots,u_m)$ of~$\R^m$, and set
$U = \begin{pmatrix}u_1&\cdots&u_m\end{pmatrix}$.

By construction, $AV = U\Sigma$ (checking on each column of~$V$), so
$A = U\Sigma\transpose{V}$.
\end{proof}

\begin{center}
\begin{tikzpicture}[>=Stealth, scale=1.6]
  % --- SVD geometry: unit circle -> ellipse ---
  % Left: unit circle in R^n
  \begin{scope}[shift={(-2.5,0)}]
    \draw[thick, blue!50!black] (0,0) circle (1);
    \draw[->, thick, red!70!black] (0,0) -- (0.707,0.707)
      node[above right] {\small $v_1$};
    \draw[->, thick, green!60!black] (0,0) -- (-0.707,0.707)
      node[above left] {\small $v_2$};
    \node[below] at (0,-1.3) {\small Unit circle in $\R^n$};
  \end{scope}

  % Arrow
  \draw[->, thick, dashed, gray] (-0.8,0) -- (0.8,0)
    node[midway, above] {\small $A$};

  % Right: ellipse in R^m
  \begin{scope}[shift={(2.5,0)}]
    \draw[thick, blue!50!black, rotate=30] (0,0) ellipse (1.5 and 0.6);
    \draw[->, thick, red!70!black] (0,0) --
      ({1.5*cos(30)},{1.5*sin(30)})
      node[above right] {\small $\sigma_1 u_1$};
    \draw[->, thick, green!60!black] (0,0) --
      ({-0.6*sin(30)},{0.6*cos(30)})
      node[above left] {\small $\sigma_2 u_2$};
    \node[below] at (0,-1.3) {\small Image ellipse in $\R^m$};
  \end{scope}
\end{tikzpicture}
\captionof{figure}{Geometric interpretation of the SVD\@.  The linear map~$A$
  sends the unit sphere to an ellipsoid.  The right singular vectors~$v_i$
  are mapped to the principal axes $\sigma_i u_i$ of the ellipsoid.}
\label{fig:svd-geometry}
\end{center}

\begin{example}{SVD of a $2\times 2$ matrix}{ex:svd-2x2}
Let $A = \begin{pmatrix}3&0\\0&-2\end{pmatrix}$.  Then
$\transpose{A}A = \begin{pmatrix}9&0\\0&4\end{pmatrix}$, with
eigenvalues $9$ and~$4$, so $\sigma_1=3$, $\sigma_2=2$.

The right singular vectors are $v_1=\vece{1}$, $v_2=\vece{2}$.
The left singular vectors are
$u_1 = \frac{1}{3}Av_1 = \vece{1}$,
$u_2 = \frac{1}{2}Av_2 = -\vece{2}$.

Thus $A = U\Sigma\transpose{V}$ with $U=\begin{psmallmatrix}1&0\\0&-1\end{psmallmatrix}$,
$\Sigma=\begin{psmallmatrix}3&0\\0&2\end{psmallmatrix}$, $V=I_2$.
\end{example}

\begin{proposition}{Properties of the SVD}{prop:svd-properties}
\index{SVD!properties}
Let $A = U\Sigma\transpose{V}$ be the SVD of $A\in\Mat_{m\times n}(\R)$
with $\rank(A)=r$.
\begin{enumerate}[label=\textup{(\roman*)}]
  \item $\rank(A)$ equals the number of nonzero singular values.
  \item $\norm{A}_2 = \sigma_1$ (operator norm) and
    $\norm{A}_F = \sqrt{\sigma_1^2+\cdots+\sigma_r^2}$ (Frobenius norm).
  \item $\im(A) = \Span(u_1,\ldots,u_r)$ and
    $\Ker(A) = \Span(v_{r+1},\ldots,v_n)$.
  \item If $A$ is square and invertible, $\inv{A} = V\inv{\Sigma}\transpose{U}$
    and the condition number is $\kappa(A) = \sigma_1/\sigma_n$.
\end{enumerate}
\end{proposition}

\begin{proof}
(i) follows from $\rank(A) = \rank(\transpose{A}A)$ (the nonzero
eigenvalues of $\transpose{A}A$ are $\sigma_1^2,\ldots,\sigma_r^2$).
(ii)--(iv) follow directly from the decomposition $A=U\Sigma\transpose{V}$.
\end{proof}


% ============================================================
\section{Applications}\label{sec:spectral-applications}
% ============================================================

% ---------- PCA ----------
\subsection{Principal Component Analysis}\label{subsec:pca}
\index{principal component analysis}\index{PCA}

Given data points $x_1,\ldots,x_N\in\R^n$ (centered so that
$\sum x_i = \vzero$), the \emph{sample covariance matrix} is
\[
  S = \frac{1}{N-1}\sum_{i=1}^{N} x_i\transpose{x_i}
  \;\in\;\Sym_n(\R).
\]
Since $S$ is symmetric and positive semidefinite, the Spectral Theorem
gives $S = P\diag(\lambda_1,\ldots,\lambda_n)\transpose{P}$ with
$\lambda_1\geq\cdots\geq\lambda_n\geq 0$.  The columns of~$P$ are the
\emph{principal components}: the first principal component~$p_1$
maximizes the variance $\transpose{p}\,S\,p$ over unit vectors~$p$,
the second maximizes the variance orthogonally to~$p_1$, and so on.

Dimensionality reduction to $k$ components amounts to projecting onto
$\Span(p_1,\ldots,p_k)$, retaining a fraction
$\frac{\lambda_1+\cdots+\lambda_k}{\lambda_1+\cdots+\lambda_n}$
of the total variance.

% ---------- Quadratic form classification ----------
\subsection{Classification of quadratic forms}%
\label{subsec:quadratic-forms}
\index{quadratic form!classification}

A \emph{quadratic form} on~$\R^n$ is a function
$Q(x) = \transpose{x}Ax$ where $A\in\Sym_n(\R)$ (since only the
symmetric part of a matrix contributes to the quadratic form).
By the Spectral Theorem, there is an orthogonal change of variables
$x = Py$ that diagonalizes~$Q$:
\[
  Q(x) = \transpose{y}\diag(\lambda_1,\ldots,\lambda_n)\,y
  = \lambda_1 y_1^2 + \cdots + \lambda_n y_n^2.
\]
The \emph{signature}\index{signature!of a quadratic form} $(p,q)$ ---
where $p$ is the number of positive~$\lambda_i$ and $q$ the number of
negative ones --- classifies the quadratic form up to congruence
(Sylvester's law of inertia\index{Sylvester's law of inertia}).

% ---------- Optimization ----------
\subsection{Second-order optimality conditions}%
\label{subsec:optimization}
\index{Hessian}\index{optimization!second-order conditions}

Let $f\colon\R^n\to\R$ be $C^2$.  At a critical point $x_0$
($\nabla f(x_0)=\vzero$), the Hessian
$H = \bigl(\frac{\partial^2 f}{\partial x_i\partial x_j}(x_0)\bigr)$
is symmetric.  By the Spectral Theorem:
\begin{itemize}
  \item $H$ positive definite $\Rightarrow$ $x_0$ is a strict local minimum.
  \item $H$ negative definite $\Rightarrow$ $x_0$ is a strict local maximum.
  \item $H$ indefinite $\Rightarrow$ $x_0$ is a saddle point.
\end{itemize}

% ---------- Vibrations ----------
\subsection{Vibrations and normal modes}\label{subsec:vibrations}
\index{vibrations}\index{normal modes}

A coupled system of $n$ harmonic oscillators (masses connected by
springs) leads to the equation $M\ddot{x} + Kx = \vzero$, where
$M,K\in\Sym_n(\R)$ are the mass and stiffness matrices (both positive
definite).  Setting $y = M^{1/2}x$ transforms this into
$\ddot{y} + M^{-1/2}KM^{-1/2}y = \vzero$.

The symmetric matrix $\tilde K = M^{-1/2}KM^{-1/2}$ has positive
eigenvalues $\omega_1^2\leq\cdots\leq\omega_n^2$; the~$\omega_i$ are
the \emph{natural frequencies} and the corresponding eigenvectors
(transformed back via $M^{-1/2}$) are the \emph{normal modes} of the
system.


% ============================================================
\section{Exercises}\label{sec:spectral-exercises}
% ============================================================

\begin{exercise}{Eigenvalues of a symmetric matrix \basic}{exo:sym-eigen}
\index{symmetric matrix!eigenvalues}
Let $A = \begin{pmatrix} 4 & 2 \\ 2 & 1 \end{pmatrix}$.
\begin{enumerate}[label=\textup{(\alph*)}]
  \item Find the eigenvalues of~$A$ and verify they are real.
  \item Find an orthonormal basis of eigenvectors.
  \item Write $A = P D \transpose{P}$ with $P$ orthogonal and $D$ diagonal.
\end{enumerate}
\end{exercise}

\begin{exercise}{Orthogonal diagonalization \basic}{exo:orth-diag}
\index{orthogonal diagonalization}
Orthogonally diagonalize
$A = \begin{pmatrix} 6 & -2 \\ -2 & 3 \end{pmatrix}$.
\end{exercise}

\begin{exercise}{A $3\times 3$ symmetric matrix \basic}{exo:3x3-sym}
Let $A = \begin{pmatrix} 1 & 0 & 1 \\ 0 & 1 & 0 \\ 1 & 0 & 1\end{pmatrix}$.
Find the eigenvalues, eigenspaces, and an orthogonal matrix~$P$ such that
$\transpose{P}AP$ is diagonal.
\end{exercise}

\begin{exercise}{Proof practice: eigenvalues of $A^2$ \basic}{exo:eigen-A2}
Let $A\in\Sym_n(\R)$ with eigenvalues $\lambda_1,\ldots,\lambda_n$.
Prove that the eigenvalues of~$A^2$ are $\lambda_1^2,\ldots,\lambda_n^2$.
\end{exercise}

\begin{exercise}{Positive definiteness tests \intermediate}{exo:pd-tests}
\index{positive definite matrix!tests}
Determine whether each matrix is positive definite, positive semidefinite,
or neither:
\begin{enumerate}[label=\textup{(\alph*)}]
  \item $A = \begin{pmatrix} 3 & 1 \\ 1 & 3 \end{pmatrix}$.
  \item $B = \begin{pmatrix} 1 & 2 \\ 2 & 1 \end{pmatrix}$.
  \item $C = \begin{pmatrix} 4 & 2 & 0 \\ 2 & 5 & 3 \\ 0 & 3 & 6
    \end{pmatrix}$.
\end{enumerate}
\end{exercise}

\begin{exercise}{Cholesky decomposition \intermediate}{exo:cholesky}
\index{Cholesky decomposition}
Compute the Cholesky decomposition $A = L\transpose{L}$ for
$A = \begin{pmatrix} 4 & 2 \\ 2 & 5 \end{pmatrix}$.
Verify your answer.
\end{exercise}

\begin{exercise}{SVD computation \intermediate}{exo:svd-compute}
\index{SVD!computation}
Compute the SVD of
$A = \begin{pmatrix} 1 & 1 \\ 0 & 1 \\ 1 & 0 \end{pmatrix}$.
\end{exercise}

\begin{exercise}{Spectral decomposition \intermediate}{exo:spectral-decomp}
\index{spectral decomposition}
Let $A = \begin{pmatrix} 3 & 1 \\ 1 & 3 \end{pmatrix}$.
Write $A$ in spectral decomposition form
$A = \lambda_1 P_1 + \lambda_2 P_2$, where $P_1,P_2$ are the
orthogonal projections onto the eigenspaces.  Verify that
$P_1+P_2 = I_2$ and $P_1P_2 = 0$.
\end{exercise}

\begin{exercise}{Normal but not self-adjoint \intermediate}{exo:normal-not-sa}
\index{normal operator}
Let $A = \begin{pmatrix} 0 & -1 \\ 1 & 0 \end{pmatrix}$.
\begin{enumerate}[label=\textup{(\alph*)}]
  \item Show that $A$ is normal.
  \item Show that $A$ has no real eigenvalues, hence is not orthogonally
    diagonalizable over~$\R$.
  \item Find the eigenvalues over~$\C$ and a unitary matrix~$U$ such
    that $\hermit{U}AU$ is diagonal.
\end{enumerate}
\end{exercise}

\begin{exercise}{Simultaneous diagonalization \challenging}{exo:simult-diag}
\index{simultaneous diagonalization}
Let $A,B\in\Sym_n(\R)$ with $AB=BA$.  Prove that $A$ and~$B$ can be
\emph{simultaneously} orthogonally diagonalized, i.e.\ there exists
$P\in\Orth(n)$ such that both $\transpose{P}AP$ and $\transpose{P}BP$
are diagonal.

\emph{Hint:} Show that each eigenspace of~$A$ is invariant under~$B$,
then apply the Spectral Theorem to the restriction of~$B$ to each
eigenspace.
\end{exercise}

\begin{exercise}{Low-rank approximation via SVD \challenging}{exo:low-rank}
\index{SVD!low-rank approximation}
Let $A\in\Mat_{m\times n}(\R)$ with SVD\
$A = \sum_{i=1}^r \sigma_i\,u_i\transpose{v_i}$.
Define $A_k = \sum_{i=1}^k \sigma_i\,u_i\transpose{v_i}$ for $k\leq r$.
\begin{enumerate}[label=\textup{(\alph*)}]
  \item Show that $\rank(A_k)=k$.
  \item Prove the Eckart--Young theorem\index{Eckart--Young theorem}:
    among all matrices of rank at most~$k$, $A_k$ minimises
    $\norm{A-B}_F$, and the minimum value is
    $\sqrt{\sigma_{k+1}^2+\cdots+\sigma_r^2}$.
\end{enumerate}
\end{exercise}

\begin{exercise}{Rayleigh quotient and eigenvalue bounds \challenging}%
{exo:rayleigh}
\index{Rayleigh quotient}
For $A\in\Sym_n(\R)$, the \emph{Rayleigh quotient}\index{Rayleigh quotient}
is $R(x) = \frac{\transpose{x}Ax}{\transpose{x}x}$ for
$x\neq\vzero$.
\begin{enumerate}[label=\textup{(\alph*)}]
  \item Prove that $\lambda_{\min}(A)\leq R(x)\leq\lambda_{\max}(A)$.
  \item Prove that $\max_{x\neq\vzero}R(x) = \lambda_{\max}(A)$
    and $\min_{x\neq\vzero}R(x) = \lambda_{\min}(A)$, and identify
    the optimisers.
  \item (\emph{Courant--Fischer minimax theorem}\index{Courant--Fischer theorem})
    Prove that the $k$-th largest eigenvalue satisfies
    \[
      \lambda_k = \max_{\dim W = k}\;\min_{x\in W,\,x\neq\vzero}
        R(x)
      = \min_{\dim W = n-k+1}\;\max_{x\in W,\,x\neq\vzero} R(x).
    \]
\end{enumerate}
\end{exercise}


% ============================================================
\section{Chapter summary}\label{sec:spectral-summary}
% ============================================================

\begin{enumerate}
  \item \textbf{Self-adjoint operators.}  An operator~$f$ is self-adjoint
    if $f^*=f$; in an orthonormal basis, this corresponds to symmetric
    ($\transpose{A}=A$) or Hermitian ($\hermit{A}=A$) matrices.

  \item \textbf{Eigenvalues are real.}  Every eigenvalue of a self-adjoint
    operator is a real number.

  \item \textbf{Orthogonal eigenvectors.}  Eigenvectors corresponding to
    distinct eigenvalues of a self-adjoint operator are orthogonal.

  \item \textbf{Spectral Theorem (real).}  Every real symmetric matrix is
    orthogonally diagonalizable: $A = P\diag(\lambda_1,\ldots,\lambda_n)\transpose{P}$
    with $P\in\Orth(n)$.

  \item \textbf{Spectral Theorem (complex).}  Every Hermitian matrix is
    unitarily diagonalizable with real eigenvalues.

  \item \textbf{Normal operators.}  An operator is normal if $ff^*=f^*f$.
    Over~$\C$, normal operators are precisely the unitarily
    diagonalizable ones.

  \item \textbf{Positive definiteness.}  $A\in\Sym_n(\R)$ is positive
    definite if and only if all eigenvalues are positive, if and only if
    all leading principal minors are positive (Sylvester's criterion).

  \item \textbf{SVD.}  Every $m\times n$ matrix admits a decomposition
    $A = U\Sigma\transpose{V}$ with $U,V$ orthogonal and $\Sigma$
    ``diagonal.''  The singular values are the square roots of the
    eigenvalues of $\transpose{A}A$.

  \item \textbf{Applications.}  The Spectral Theorem underlies PCA,
    quadratic form classification, second-order optimality conditions,
    and the analysis of coupled vibrations.
\end{enumerate}
